\section{Theoretical Analysis}\label{sec:theory}

This section answers two design questions: when SOMA may accept the session-local surrogate, and when its warm-start overhead is amortized. \rev{Both answers assume that later turns stay tied to the warm-start state, though not that they are easier. The results describe when SOMA can operate and claim no advantage over other mining strategies. Appendix~\ref{sec:app_proof} has the proofs and a candidate bound.}

\subsection{Switching Criterion}\label{sec:warm-start}

At turn $t$, with shared context $S_t=\mathcal D_{t-1}\oplus q_t$, we measure the discrepancy of the adapted surrogate with the bounded score $\mathsf{Gap}_{\Theta}(S_t)\in[0,1]$ of Section~\ref{sec:prelim}, where $\Theta$ is the current adapter. \rev{Since consecutive turns are dependent, we model per-turn statistics as an $m$-dependent sequence and use the effective sample sizes $W_{\mathrm{eff}}=W/(m+1)$ for the warm-start window and $|B|_{\mathrm{eff}}=|B|/(m+1)$ for an acceptance batch $B$. Lemma~\ref{lem:warmstart} shows that, with probability at least $1-\delta$, a bounded warm-start statistic lies within $\eta_W(\delta)=\sqrt{\log(2/\delta)/(2W_{\mathrm{eff}})}$ of its mean, so strongly dependent early turns may need a longer warm start before the switch can be assessed.}

\rev{The fidelity gate evaluates the adapted surrogate on an acceptance batch $B$ through its average gap $\widehat{\mathsf G}_B=\frac{1}{|B|}\sum_{S_j\in B}\mathsf{Gap}_{\Theta}(S_j)$, and the surrogate passes when $\widehat{\mathsf G}_B$ is below a threshold $\varepsilon$.}

\begin{proposition}[Acceptance Bound for Switching]\label{thm:detection}
\rev{Assume that warm-start and later states share the distribution $\mathcal Q_k$ (the locality premise of Section~\ref{sec:prelim}), let $B$ be drawn from it and be disjoint from the turns used to fit $\Theta$, and let $\gamma_B(\delta)=\sqrt{\log(1/\delta)/(2|B|_{\mathrm{eff}})}$. With probability at least $1-\delta$, $\mathbb E_{S\sim\mathcal Q_k}[\mathsf{Gap}_{\Theta}(S)]\le\widehat{\mathsf G}_B+\gamma_B(\delta)$. Hence $\widehat{\mathsf G}_B\le\varepsilon-\gamma_B(\delta)$ certifies $\mathbb E[\mathsf{Gap}_{\Theta}]\le\varepsilon$.}
\end{proposition}

\rev{A union bound combines the two results (Corollary~\ref{cor:switch}). An abrupt topic change typically fails the locality gate and triggers an immediate rollback to $F$. The guarantee has two practical limits. It needs a held-out acceptance batch, but our implementation reuses the warm-start turns for adaptation and acceptance, so in the reported experiments the gate works as a heuristic check. The bound is also loose. With $\delta=0.05$ and $|B|=24$, $\gamma_B\approx0.25$, which exceeds the thresholds $\varepsilon\in[0.05,0.12]$ themselves, so we set these thresholds empirically (Appendix~\ref{sec:app_proof}).}

\subsection{Net Efficiency Gain}\label{sec:theory-cost}

Because SOMA pays its warm-start overhead once, its efficiency must be judged per session. Let $T$ be the number of turns, $W$ the switch point, and $c_F(\cdot)$ and $c_G(\cdot)$ the per-turn costs (latency, tokens, or price) of $F$ and of $G$ on its compressed input $\widetilde S_t$; the net gain over serving every turn with $F$ is
\begin{equation}
\label{eq:net-gain}
\Delta C(T;W)=\sum_{t=W+1}^{T}\Big[c_F(\mathcal D_{t-1}\oplus q_t)-c_G(\widetilde S_t)\Big]
-C_{\mathrm{probe}}(W)-C_{\mathrm{LoRA}}(W)-C_{\mathrm{rb}}(T;W),
\end{equation}
where $C_{\mathrm{probe}}$ and $C_{\mathrm{LoRA}}$ are the mining and adaptation costs and $C_{\mathrm{rb}}$ covers drift checks and rollback; SOMA pays off only when the accumulated post-switch savings exceed this one-time overhead. \rev{This is consistent with Figure~\ref{fig:operating}(a), where sessions of 1--4 turns do not benefit and locally coherent sessions of 5 or more turns recover the initialization cost.}
